\begin{table}[H]
  \begin{center}
    \setlength\extrarowheight{-6pt}
    \begin{tabular}{c|ccccccccc}
      $n/k$ & 0 & 1 & 2 & 3 & 4 & 5 & 6 & 7 & 8 \\ [3px]
      \hline
      0 & 1 &   &   &   &   &   &   &   &   \\
      1 & 2 & 1 &   &   &   &   &   &   &   \\
      2 & 4 & 4 & 1 &   &   &   &   &   &   \\
      3 & 8 & $\tfrac{49}{4}$ & 6 & 1 &   &   &   &   &   \\
      4 & 16 & 34 & 25 & 8 & 1 &   &   &   &   \\
      5 & 32 & $\tfrac{1441}{16}$ & 90 & $\tfrac{85}{2}$ & 10 & 1 &   &   &   \\
      6 & 64 & $\tfrac{931}{4}$ & 301 & 190 & 65 & 12 & 1 &   &   \\
      7 & 128 & $\tfrac{37969}{64}$ & 966 & $\tfrac{12411}{16}$ & 350 & $\tfrac{371}{4}$ & 14 & 1 &   \\
      8 & 256 & $\tfrac{6001}{4}$ & 3025 & 3003 & 1701 & 588 & 126 & 16 & 1 \\
    \end{tabular}
  \end{center}
  \caption{Values of generalized central factorial numbers $T_2(n,k)$.
    See also the sequences \href{https://oeis.org/A000000}{\texttt{A000000}},
    and \href{https://oeis.org/A000000}{\texttt{A000000}} in~\cite{sloane2003line}.
  }
\end{table}
